\documentclass[12pt]{article}
\usepackage{fullpage}
\usepackage{amsmath,amssymb,amsthm}

\newtheorem{lemma}{Lemma}
\newtheorem{proposition}{Proposition}
\newtheorem{theorem}{Theorem}

\newcommand{\Dil}{\operatorname{Dil}}
\newcommand{\Vol}{\operatorname{Vol}}
\newcommand{\Fill}{\operatorname{Fill}}
\newcommand{\Mass}{\operatorname{Mass}}
\newcommand{\spt}{\operatorname{spt}}
\newcommand{\RR}{\mathbb R}

\begin{document}

\begin{center}
{\bf Partial proof and a sharpened random-puncture reduction}
\end{center}

\section*{Problem statement and interpretation}

Rectangles are oriented Euclidean boxes.  A map
\(f:(R,\partial R)\to(S,\partial S)\) has degree \(1\) if
\(f_*[R,\partial R]=[S,\partial S]\) in relative homology.  I write the
small universal constant in the problem as \(\kappa\), to avoid confusing it
with \(k\)-dilation.

The argument below proves the standard first alternative and reduces the
remaining case to a precise family-level deformation estimate, Lemma
\ref{lem:random-puncture}.  The new reduction is stronger than the false
averaged estimate used in an earlier draft: the uncontrolled part is charged to
the actual minimizing fillings and therefore has the correct self-absorbing
coefficient.  The proof is still incomplete exactly because Lemma
\ref{lem:random-puncture} is not proved here.  All implicit constants are
positive dimensional constants.  Strict inequalities in the problem would follow
from non-strict estimates after decreasing \(\kappa\).

\section*{Standard tools}

We use integral Lipschitz chains/currents.  Let \(T=[R,\partial R]\).  Since
\(f\) has relative degree \(1\), the relative current
\(f_\#T-[S,\partial S]\) has zero boundary and zero relative homology class; by
the constancy theorem it is zero in \(S\) modulo currents supported in
\(\partial S\).  Thus, if \(\pi_{ij}\) denotes projection to two target
coordinates and \(F=\pi_{ij}\circ f\), the naturality of slicing gives
\[
        f_\#\langle T,F,y\rangle
        =\langle [S,\partial S],\pi_{ij},y\rangle
\]
for a.e. regular value \(y\), as relative currents.  We use Guth's open-rectangle
isoperimetric estimates for closed rectangles by applying them in an
\(\varepsilon\)-collared interior and letting \(\varepsilon\downarrow0\).

If \(\Dil_2(f)\le1\), then \(\Dil_j(f)\le1\) for \(j=3,4\).  Indeed, if
\(\sigma_1\ge\cdots\ge\sigma_4\) are the singular values of \(df\), then
\(\sigma_1\sigma_2\le1\), hence \(\sigma_i\le1\) for \(i\ge2\), and
\(\sigma_1\cdots\sigma_j\le1\) for \(j\ge2\).

We use two estimates of Guth.  First, the \(k=2,n=4\) small-volume part of the
rectangular relative isoperimetric profile theorem \cite[Theorem 3]{GuthExpanding}
says that if \(Z\) is a relative integral \(2\)-cycle in \(R\),
\(A=\Mass Z\le c_I R_1R_2\), then
\[
        \Fill_R(Z)\le C_I
        \max\left\{ A^{3/2},\,\frac{A^2}{R_1}\right\}.        \tag{1}
\]
In particular, in this range \(\Fill_R(Z)\le C_I R_1R_2^2\).  Second, Guth's
degree estimate \cite[Estimate 1]{GuthExpanding}, with \(j=1,k=2,l=4\), gives
\[
        R_1^3R_2R_3R_4\ge c_G S_1^3S_2S_3S_4                  \tag{2}
\]
for every degree-one map of pairs with \(\Dil_2(f)\le1\).

\section*{The first alternative}

\begin{lemma}[central target filling]\label{lem:target-fill}
Let
\[
 P_y=[0,S_1]\times[0,S_2]\times\{y_3\}\times\{y_4\},
\]
where \(y_3\in[S_3/3,2S_3/3]\) and
\(y_4\in[S_4/3,2S_4/3]\).  Every relative \(3\)-chain \(Y\) in \(S\) with
\(\partial Y=P_y\) in \(I_2(S,\partial S)\) obeys
\[
        \Mass Y\ge c\,S_1S_2S_3 .
\]
\end{lemma}

\begin{proof}
Choose a smooth \(1+o(1)\)-Lipschitz function \(\psi\) on the
\((x_3,x_4)\)-rectangle which is zero on the boundary and satisfies
\(\psi(y_3,y_4)\ge cS_3-o(1)\).  Apply Stokes' theorem to
\(\omega=dx_1\wedge dx_2\wedge d\psi\).  The comass of \(\omega\) is at most
\(1+o(1)\), so
\[
 \Mass Y\ge(1+o(1))^{-1}\left|\int_Y\omega\right|
  =(1+o(1))^{-1}\left|\int_{\partial Y}\psi\,dx_1\wedge dx_2\right|.
\]
The part of \(\partial Y\) lying in \(\partial S\) contributes zero: on the
\(x_1\)- and \(x_2\)-faces \(dx_1\wedge dx_2=0\), while on the \(x_3\)- and
\(x_4\)-faces \(\psi=0\).  Letting \(o(1)\to0\) proves the claim.
\end{proof}

\begin{proposition}[first alternative]\label{prop:first}
There are constants \(c_1,C_1>0\) such that if
\(f:(R,\partial R)\to(S,\partial S)\) has degree \(1\), \(\Dil_2(f)\le1\), and
\[
        C_1R_1R_2^2\le S_1S_2S_3 ,
\]
then \(R_3R_4\ge c_1S_3S_4\).
\end{proposition}

\begin{proof}
Let \(F=(x_3,x_4)\circ f\), and set
\[
Q=[S_3/3,2S_3/3]\times[S_4/3,2S_4/3].
\]
For a.e. \(y\in Q\), put \(Z_y=\langle [R,\partial R],F,y\rangle\).  Slicing
naturality gives \(f_\# Z_y=P_y\) as relative currents, up to orientation.
Thus any relative filling of \(Z_y\) in \(R\) pushes forward to a relative
filling of \(P_y\) in \(S\).  Since \(\Dil_3(f)\le1\), Lemma
\ref{lem:target-fill} implies
\[
        \Fill_R(Z_y)\ge cS_1S_2S_3 .                         \tag{3}
\]
Choose \(C_1\) large compared with \(C_I/c\).  If
\(C_1R_1R_2^2\le S_1S_2S_3\), then (1) and (3) force
\(\Mass Z_y\ge c_I R_1R_2\) for a.e. \(y\in Q\).  Coarea and \(J_2F\le1\) give
\[
 \Vol(R)\ge\int_Q\Mass Z_y\,dy
      \ge c R_1R_2S_3S_4 .
\]
Since \(\Vol(R)=R_1R_2R_3R_4\), the proposition follows.
\end{proof}

\section*{Consequences of the same slicing}

For the same central fibers,
\[
 \Mass Z_y\ge c\min\left\{R_1R_2,\,
          (S_1S_2S_3)^{2/3},\,
          (R_1S_1S_2S_3)^{1/2}\right\}                       \tag{4}
\]
for a.e. \(y\in Q\).  Indeed, if \(A=\Mass Z_y<c_I R_1R_2\), then (1) and
(3) imply
\[
 cS_1S_2S_3\le C_I\max\{A^{3/2},A^2/R_1\}.
\]
After integration over \(Q\),
\[
 \Vol(R)\ge cS_3S_4\min\left\{R_1R_2,\,
          (S_1S_2S_3)^{2/3},\,
          (R_1S_1S_2S_3)^{1/2}\right\}.                      \tag{5}
\]
Also, (2) and \(R_1\le\kappa S_1\) give
\[
        \Vol(R)\ge c_G\kappa^{-2}S_1S_2S_3S_4 .              \tag{6}
\]
Thus (6) implies the desired second alternative whenever
\(S_3/S_2\le c\,\kappa^{-6}\).

It remains only to discuss the case in which the first desired alternative
fails.  Taking \(\kappa<\min\{c_1,1\}\), Proposition \ref{prop:first} gives
\[
        R_1R_2^2\gtrsim S_1S_2S_3 .                          \tag{7}
\]
Put
\[
        \alpha=\frac{R_1}{S_1},\qquad
        T=S_1S_2^{1/2}S_3^{3/2}S_4 .
\]
Then the three terms in (5) are all at least
\(c(R_1S_1S_2S_3)^{1/2}\), and hence
\[
        \Vol(R)\ge c\,\alpha^{1/2}T .                         \tag{8}
\]
The monomial estimate gives the sharper form
\[
        \Vol(R)\ge c\,\alpha^{-2}S_1S_2S_3S_4 .                \tag{9}
\]
Consequently, if the second desired alternative also fails, then
\[
        \alpha\lesssim \kappa^2,\qquad
        \frac{S_3}{S_2}\gtrsim \alpha^{-4}\kappa^{-2}.          \tag{10}
\]
Thus only an extreme high-aspect-ratio regime remains.

\section*{Weighted coarea}

Let \(F=(f_3,f_4)\), \(G=(f_1,f_2)\), and, at rank-two points of \(F\), set
\[
        \lambda(x)=\left\|dG|_{\ker dF_x}\right\|,
\]
with \(\lambda=0\) on the rank-deficient set.  For
\(E_y=\int \lambda^2\,d\|Z_y\|\) and \(I=\int_Q E_y\,dy\), one has
\[
        \Vol(R)\ge I,                                      \tag{11}
\]
and, for a.e. \(y\in Q\),
\[
        E_y\ge S_1S_2 .                                    \tag{12}
\]
Indeed, if \(T_x=\ker dF_x\), \(N_x=T_x^\perp\), and \(v\in T_x\) is unit with
\(|dG(v)|=\lambda\), then for every unit \(n\in N_x\),
\[
        |df(v)\wedge df(n)|\ge \lambda |dF(n)|.
\]
Since \(\Dil_2(f)\le1\), this implies \(\lambda^2J_2F\le1\), and (11) follows
from coarea.  For (12), slicing naturality gives
\(G_\#Z_y=[0,S_1]\times[0,S_2]\) as a relative current; the area formula and
\(\lambda^2\ge J_2(G|_{Z_y})\) give the claim.  Alone this recovers only
\(\Vol(R)\gtrsim S_1S_2S_3S_4\).

\section*{Projection and the random-puncture reduction}

Let \(q=|Q|\sim S_3S_4\), \(F_0=S_1S_2S_3\), and
\[
        A=\int_Q\Fill_R(Z_y)\,dy .
\]
Choose measurable almost-minimizing fillings \(B_y\) of \(Z_y\), so that
\(\partial B_y=Z_y\) relatively and
\(\int_Q\Mass B_y\,dy\le2A\).  For \(u=(x_1,x_2)\), write
\(V_u=\{u\}\times[0,R_3]\times[0,R_4]\).  Given \(z\in Q\), define
\[
        \mathcal B_z=\{x\in R:\ z\in F(V_{\pi_{12}x})\}.
\]
This is the set of points whose vertical rectangle sees the puncture \(z\).
Fubini and the area bound \(J_2(F|_{V_u})\le1\) give
\begin{equation}
\begin{aligned}
 \frac{1}{q}\int_Q\int_Q \Mass(B_y\llcorner\mathcal B_z)\,dy\,dz
 &\le \frac{1}{q}\int_Q\int_R
      \mathcal L^2(F(V_{\pi_{12}x})\cap Q)\,d\|B_y\|(x)\,dy      \\
 &\le \frac{R_3R_4}{q}\int_Q\Mass B_y\,dy
 \le 2\,\frac{R_3R_4}{S_3S_4}\,A .
\end{aligned}                                                   \tag{13}
\end{equation}
Thus a random puncture is bad for only the correct self-absorbing fraction of the
actual fillings, not merely for an unweighted fraction of parameter space.

The remaining missing ingredient can now be isolated.

\begin{lemma}[random-puncture deformation]\label{lem:random-puncture}
There is a dimensional constant \(C_A\) such that for every \(z\in Q\),
\[
 A\le C_A\left(R_1I+\frac{I^2}{S_1q}\right)
     +C_A\int_Q \Mass(B_y\llcorner\mathcal B_z)\,dy .          \tag{14}
\]
\end{lemma}

This lemma is the sought two-parameter Guth/Federer--Fleming tightening.  The
puncture is used as a stopping obstruction: on \(R\setminus\mathcal B_z\), every
vertical rectangle \(V_u\) has \(F(V_u)\subset Q\setminus\{z\}\), so the vertical
part of the two-parameter slice family is null in the punctured parameter
rectangle.  The expected controlled costs are \(R_1I\), from pushing in the short
\(x_1\)-direction weighted by the mixed energy, and \(I^2/(S_1q)\), from capping
the resulting \(f_1\)-level seams by a one-dimensional coarea optimization.  I do
not have a rigorous proof of this deformation lemma.

\begin{proposition}[the deformation lemma would finish the proof]\label{prop:selfabsorb}
If Lemma \ref{lem:random-puncture} holds, then the theorem in the problem follows.
\end{proposition}

\begin{proof}
Average (14) in \(z\), use (13), and choose a good puncture.  This gives
\[
 A\le C_A\left(R_1I+\frac{I^2}{S_1q}\right)
       +C_A'\frac{R_3R_4}{S_3S_4}A .                         \tag{15}
\]
By (3), \(A\ge cF_0q\).  If the first desired alternative fails and
\(\kappa<1/(2C_A')\), the last term in (15) is absorbed, so
\[
        F_0q\lesssim R_1I+\frac{I^2}{S_1q}.                  \tag{16}
\]
If \(I\le \eta qS_1(S_2S_3)^{1/2}\), then
\[
        R_1I\le \eta\kappa F_0q,\qquad
        \frac{I^2}{S_1q}\le \eta^2F_0q ,
\]
using \(R_1\le\kappa S_1\) and \(S_1\le S_2\le S_3\).  Choosing \(\eta\) and then
\(\kappa\) small contradicts (16).  Hence
\[
        I\ge c\,q\,S_1(S_2S_3)^{1/2}.
\]
Finally (11) yields
\[
 \Vol(R)\ge I
   \ge c S_3S_4\,S_1(S_2S_3)^{1/2}
   =cS_1S_2^{1/2}S_3^{3/2}S_4 ,
\]
which is the second desired alternative after decreasing \(\kappa\).
\end{proof}

\section*{The false non-absorbing estimate}

It is important that (15) contains \(A\) in the last term.  The stronger-looking
non-absorbing estimate with
\((R_3R_4/S_3S_4)F_0q\) in place of that term is false.  Let
\(S=(1,1,L,L)\), \(R=(1,M,L,L)\), \(1<M<L\), and
\[
        f(u_1,u_2,u_3,u_4)=(u_1,u_2/M,u_3,u_4).
\]
Then \(\Dil_2(f)=1\), the map has degree \(1\), \(I\sim ML^2\), and
\(q\sim L^2\).  The central slices are
\([0,1]\times[0,M]\times\{y_3\}\times\{y_4\}\), and the calibration
\(du_1\wedge du_2\wedge d\psi(u_3,u_4)\) gives
\(\Fill_R(Z_y)\gtrsim ML\); hence \(A\gtrsim ML^3\).  The three non-absorbing
terms would be \(ML^2\), \(M^2L^2\), and \(L^3\), which are too small when
\(M=L^{1/2}\).

\section*{Remaining open issues}

The proof is incomplete because Lemma \ref{lem:random-puncture}, equivalently
the self-absorbing estimate (15), has not been proved.

\begin{itemize}
\item Mathematically, the missing point is a two-parameter
Federer--Fleming/Guth tightening estimate for the slice family
\(Z_y=\langle [R,\partial R],(f_3,f_4),y\rangle\), using the mixed pointwise
bound \(\lambda^2J_2F\le1\).
\item The gap appears exactly in Lemma \ref{lem:random-puncture}.  The random
puncture formulation fixes the earlier false independence step by charging the
bad part to the actual fillings \(B_y\), as shown in (13).
\item What has been proved without this lemma: the first alternative,
the lower bounds (5), (8), and (9), and the extreme reduction (10).  Thus any
counterexample to the desired theorem would have to lie in a very anisotropic
regime with \(R_1/S_1\ll\kappa\) and \(S_3/S_2\) correspondingly enormous.
\item What is needed: a rigorous proof of the deformation in
\(R\setminus\mathcal B_z\), most likely by discretizing \(Q\), applying Guth's
tightening skeleton-by-skeleton to the family of slices, and using the puncture
to avoid paying the long \(R_3,R_4\) vertical cost except on the bad mass in
(14).
\end{itemize}

\begin{thebibliography}{9}
\bibitem{GuthExpanding}
L. Guth, \emph{Area-expanding embeddings of rectangles}, arXiv:0710.0403,
especially Theorem 3 and Estimate 1.

\bibitem{GuthThesis}
L. Guth, \emph{Area-contracting maps between rectangles}, Ph.D. thesis, MIT,
2005.

\bibitem{Federer}
H. Federer, \emph{Geometric Measure Theory}, Springer, 1969.
\end{thebibliography}

\end{document}
